\documentclass{article}
\usepackage{graphicx}
\usepackage{listings}

\lstset { %
    language=C++
}

\title{WG01: C++ for Kernel Development}
\author{Amlal El Mahrouss\\NeKernel.org}
\date{\today}

\begin{document}

\maketitle

\section{Abstract}
{
Many and most kernels have been shipped using the C programming language.\\
And some of them like EKA2 uses the C++ programming language. Although notoriously difficult, one may still adapt to those constraints in order to deliver one such operating system kernel. \\
That is the reason that most production-grade kernels (Linux, XNU, and NT) are mostly written in C. With a higher-level subset in C++. \\
However, when correctly applying C++ principles to kernel development, one makes the development much more easier to pull off.
}

\section{The Three Principles of Kernel C++}
\subsection{Part One: The C++ Runtime}
{
The problem mostly lies in the C++ runtime. Which assumes an existing host. A host is the contrary of a freestanding target, that is a program which expects a runtime to be present and linked to the program. \\
A C++ Kernel may instead make use of compile-time features of C++ alongside a tiny C++ runtime to make sure that no issues arise because of this host/freestanding difference.
}

\subsection{Part Two: Constexpr and Friends}
{
One may avoid v-tables and runtime dependent features as much as possible. \\
While focusing instead on meta-programming and compile-time features offered by C++. For example one may use templates to implement a scheduling policy algorithm. \\
One example of such implementation may be: \\

(Keep in mind that the code has not been tested on production systems, use it at your own risk!)

\begin{lstlisting}
/// Reference Implementation: https://github.com
/// /nekernel-org/nekernel/blob/stable/src/kernel/
/// KernelKit/CoreProcessScheduler.h#L51-L76
/// AND: https://github.com/nekernel-org/nekernel/blob/stable/src/kernel
/// /KernelKit/CoreProcessScheduler.h#L78-L105

struct FileTree final {
    static constexpr bool is_virtual_memory = false;
    static constexpr bool is_memory = false;
    static constexpr bool is_file = true;
    /// ...
};

struct MemoryTree final {
    static constexpr bool is_virtual_memory = false;
    static constexpr bool is_memory = true;
    static constexpr bool is_file = false;
    /// ...
};
\end{lstlisting}

As you can see these two structure 
leverages the `constexpr' keyword to make sure
no bugs or panic occurs at runtime\\because of a misuse of a system resource. \\

Which is why the constexpr keyword is very powerful here, we avoid the many pitfalls of writing (and hoping) that the C version will be well-thought enough so that we can catch such bugs later.

\subsection{Bonus: Using Concepts and the DDK.}
{
This bonus subsection intends to show how properly applied C++\\
Can solve issues related to driver validation.
The following example shows how powerful C++ can be when combining it with Device Driver development as well.

\begin{lstlisting}
/// Reference implementation:
/// https://github.com/nekernel-org
/// /nekernel/blob/stable/src/libDDK/DriverKit
/// /c%2B%2B/driver_base.h
/// @brief This concept requires the driver 
/// to be IDriverBase compliant.

template <typename T>
concept IsValidDriver = requires(T a) {
  { a.IsActive() && a.Type() > 0 };
};

/// @brief Consteval helper to detect
/// whether a template is truly based on IDriverBase.
/// @note This helper is consteval only.
template<IsValidDriver T>
consteval void ce_ddk_is_valid(T) {}
\end{lstlisting}
}
}

\subsection{Part Three: Memory and C++ Classes}
{
This last part treats about the final and most important part of this paper so far. Memory. As you may already have known, the C++ language dress a class lookup system (also called v-table) in order to refer to a base method in case if the instance has it missing.

\begin{lstlisting}
/// Link: https://godbolt.org/z/q3Wceannr
/// /std:c++20 /Wall

#include <iostream>

class A
{
public:
    explicit A() = default;
    virtual ~A() = default;

    virtual void doImpl() 
    {
        std::cout << "doImpl\n";
    }
};

class B : public A
{
public:
    explicit B() = default;
    ~B() override = default;
};

int main() {
    B* callImpl = new B();
    callImpl->doImpl();
}
\end{lstlisting}

\subsection{Addressing V-Tables in a C++ Kernel.}

Now the problem with kernel development is that we want to avoid such feature as much as possible, and we'd do that by following the Prong on Inheritance:

\subsection{The Three Prongs on Inheritance:}

Consider the following questions:

\begin{itemize}
\item[A:] Is this a protocol/concept that can be extended to other similar protocols/concepts?
\item[B:] Can I do this without too much trade-off costs?
\item[C:] Can I do this without using V-Tables?
\end{itemize}
}

When 2/3 of those questions fail, you should consider finding another solution to your problem, as it surely has an equivalent without V-Tables.

\section{Final Words and Conclusion}
{
We can now conclude that C++ in Kernel Development is indeed possible under strict conditions of C++ Development. \\
Breaking those conditions would lead to system quirks and instability.
A reference implementation of this paper exists, It's called NeKernel.org, available under the same internet address. \\
I am looking forward to any questions or inquiries at: amlal@nekernel.org \\\\Yours truly, \\
Amlal El Mahrouss
}

\section{References}
{
\begin{itemize}
    \item[EKA2:] https://en.wikipedia.org/wiki/EKA2
    \item[NeKernel.org:] https://nekernel.org
\end{itemize}
}

\end{document}
